\begin{figure}[htbp]
\centering
\begin{tikzpicture}
\begin{polaraxis}[
    name=radar,
    width=0.78\textwidth,
    height=0.78\textwidth,
    xtick={0,60,120,180,240,300},
    xticklabels={
        {Генерация\\из документа},
        Методология,
        {Русский\\язык},
        {Локальное\\развёртывание},
        {Проверка\\фактов},
        {Автоматическое\\обновление}
    },
    xticklabel style={font=\footnotesize, anchor=center, align=center, inner sep=8pt},
    ymin=0, ymax=1.2,
    ytick={0.25,0.5,0.75,1},
    yticklabels={},
    y grid style={gray!30},
    x grid style={gray!50, thick},
    grid=both,
    legend style={
        at={(0.5,-0.1)},
        anchor=north,
        font=\footnotesize,
        cells={anchor=west},
        legend columns=3,
        column sep=6pt,
    },
]

% ИИ-методист (рисуем первым — фон)
\addplot[teal!70!black, very thick, mark=*, mark size=2.5pt, fill=teal!12!white]
    coordinates {(0,1)(60,1)(120,1)(180,1)(240,1)(300,1)(360,1)};
\addlegendentry{ИИ-методист}

% Coursebox
\addplot[blue!70!black, thick, mark=square*, mark size=2.5pt]
    coordinates {(0,1)(60,0)(120,0.5)(180,0)(240,0)(300,0.5)(360,1)};
\addlegendentry{Coursebox}

% Articulate AI
\addplot[red!70!black, thick, mark=triangle*, mark size=2.5pt]
    coordinates {(0,0.5)(60,0)(120,0.5)(180,0)(240,0)(300,0)(360,0.5)};
\addlegendentry{Articulate AI}

% iSpring AI
\addplot[orange!80!black, thick, mark=diamond*, mark size=2.5pt]
    coordinates {(0,0.5)(60,0)(120,1)(180,0)(240,0)(300,0)(360,0.5)};
\addlegendentry{iSpring AI}

% Synthesia
\addplot[purple!70!black, thick, mark=pentagon*, mark size=2.5pt]
    coordinates {(0,0)(60,0)(120,0.5)(180,0)(240,0)(300,0)(360,0)};
\addlegendentry{Synthesia}

% Подписи шкалы вдоль одного луча (90° = вверх-вправо, ось "Методология")
\node[font=\tiny, fill=white, inner sep=1pt, text=gray!70!black] at (axis cs:30,0.25) {0{,}25};
\node[font=\tiny, fill=white, inner sep=1pt, text=gray!70!black] at (axis cs:30,0.5) {0{,}5};
\node[font=\tiny, fill=white, inner sep=1pt, text=gray!70!black] at (axis cs:30,0.75) {0{,}75};
\node[font=\tiny, fill=white, inner sep=1pt, text=gray!70!black] at (axis cs:30,1.0) {1{,}0};

\end{polaraxis}
\end{tikzpicture}
\caption{Радарная диаграмма сравнения решений}
\label{fig:competitor-radar}
\end{figure}
